\begin{proposition}[5.9.11]
\label{IV.2.s5.9.11-fr}
Bajo las hipótesis de \textup{(5.9.8)}, los $\mathcal O_X$-módulos
$\operatorname{Ker}(\rho_{X/Z})$ y $\operatorname{Coker}(\rho_{X/Z})$ tienen
su soporte contenido en $Z$, y el $\mathcal O_X$-módulo
$\mathcal H^0_{X/Z}(\mathcal F)$ es $Z$-cerrado. Además, si
$u:\mathcal F\to\mathcal F'$ es un homomorfismo de
$\mathcal O_X$-módulos tal que $\mathcal F'$ sea $Z$-cerrado, $u$ se factoriza
de una única manera en
$\mathcal F\xrightarrow{\rho_{X/Z}}\mathcal H^0_{X/Z}(\mathcal F)
\xrightarrow{v}\mathcal F'$.
Si, además, los soportes de $\operatorname{Ker}(u)$ y de
$\operatorname{Coker}(u)$ están contenidos en $Z$, $v$ es un isomorfismo.
\end{proposition}

La primera afirmación significa que, para todo $x\in X\setmin Z$, se tiene
\[
\operatorname{Ker}(\rho_{X/Z})_x
=\operatorname{Coker}(\rho_{X/Z})_x=0,
\]
es decir, que $(\rho_{X/Z})_x$ es biyectivo, o también que $\mathcal F$ es
$Z$-puro en $x$, como se ha señalado más arriba (5.9.9).

Para mostrar que $\mathcal H^0_{X/Z}(\mathcal F)$ es $Z$-cerrado, se trata de
ver que, para un abierto noetheriano $V$ de $X$,
$\varinjlim_\beta
\Gamma(V\cap U_\beta,\mathcal H^0_{X/Z}(\mathcal F))$
es igual a $\Gamma(V,\mathcal H^0_{X/Z}(\mathcal F))$; pero, por definición,
$\Gamma(V\cap U_\beta,\mathcal H^0_{X/Z}(\mathcal F))
=\varinjlim_\alpha\Gamma(V\cap U_\alpha\cap U_\beta,\mathcal F)$.
Ahora bien, la familia doble $(U_\alpha\cap U_\beta)$ es filtrante decreciente,
y nuestra afirmación resulta de (5.9.1) y del teorema del doble límite
inductivo.

\oldpage[IV-2]{113}

Pasemos a la segunda parte de la proposición. La existencia y la unicidad de
$v$ resultan de los hechos siguientes. Para todo $\alpha$, $(i_\alpha)_*(u)$
es el único homomorfismo que hace conmutativo el diagrama
\[
\xymatrix{
\mathcal F\ar[r]^u\ar[d] & \mathcal F'\ar[d]\\
(i_\alpha)_*(\mathcal F|U_\alpha)
\ar[r]_{(i_\alpha)_*(u)} &
(i_\alpha)_*(\mathcal F'|U_\alpha)
}
\]
Además, existe un único homomorfismo $w$ que hace conmutativos todos los
diagramas
\[
\xymatrix{
(i_\alpha)_*(\mathcal F|U_\alpha)
\ar[r]^{(i_\alpha)_*(u)}\ar[d] &
(i_\alpha)_*(\mathcal F'|U_\alpha)\ar[d]\\
\mathcal H^0_{X/Z}(\mathcal F)\ar[r]_w &
\mathcal H^0_{X/Z}(\mathcal F')
}
\]
Finalmente, $\mathcal F'$ y $\mathcal H^0_{X/Z}(\mathcal F')$ se identifican
canónicamente por hipótesis.

Queda por ver que, si los soportes de $\operatorname{Ker}(u)$ y de
$\operatorname{Coker}(u)$ están contenidos en $Z$, $v$ es un isomorfismo.
Basta ver que, para todo abierto noetheriano $V$, el homomorfismo
correspondiente
$\Gamma(V,\mathcal H^0_{X/Z}(\mathcal F))
\longrightarrow\Gamma(V,\mathcal F')$
es entonces un isomorfismo. Ahora bien, si una sección
$t\in\Gamma(V,\mathcal H^0_{X/Z}(\mathcal F))$ tiene imagen $0$ en
$\Gamma(V,\mathcal F')$, observemos que, para algún índice $\alpha$, se tiene
$t\in\Gamma(V\cap U_\alpha,\mathcal F)$ y, en virtud de la hipótesis sobre
$u$, se tiene $t_y=0$ para todo $y\in V\cap(X\setmin Z)$; existe, por consiguiente,
un abierto que contiene $V\cap(X\setmin Z)$ tal que la restricción de $t$ a este
abierto es nula, por tanto, por definición, $t$ es el elemento $0$ de
$\Gamma(V,\mathcal H^0_{X/Z}(\mathcal F))$.

Probemos ahora que toda sección $s'\in\Gamma(V,\mathcal F')$ es la imagen de
una sección de $\mathcal H^0_{X/Z}(\mathcal F)$ sobre $V$. Por hipótesis,
para todo $x\in V\cap(X\setmin Z)$ existe una sección $s^{(x)}$ de $\mathcal F$
sobre una vecindad abierta $W^{(x)}$ de $x$ en $X$, cuya imagen por $u$ es
$s'|W^{(x)}$; $s'|W^{(x)}$ es, por tanto, también la imagen por $v$ de la
sección $t^{(x)}$ de $\mathcal H^0_{X/Z}(\mathcal F)$, imagen canónica de
$s^{(x)}$. Además, como se ha visto que $v$ es inyectivo, las restricciones
de $t^{(x)}$ y $t^{(x')}$ a $W^{(x)}\cap W^{(x')}$ son idénticas para dos
puntos cualesquiera $x$, $x'$ de $V\cap(X\setmin Z)$; los $t^{(x)}$ son, por
consiguiente, restricciones de una misma sección $t$ de
$\mathcal H^0_{X/Z}(\mathcal F)$ sobre una vecindad abierta $U$ de
$(X\setmin Z)\cap V$. Pero, como $\mathcal H^0_{X/Z}(\mathcal F)$ es $Z$-cerrado,
$t$ se prolonga de una única manera a una sección de
$\mathcal H^0_{X/Z}(\mathcal F)$ sobre $V$, cuya imagen por $v$ tiene la
misma restricción que $s'$ en $U$, y coincide, por tanto, con $s'$ por la misma
razón.

Se dice que $\mathcal H^0_{X/Z}(\mathcal F)$ es la $Z$-\emph{clausura} de
$\mathcal F$.

\oldpage[IV-2]{114}
